% SGA 3, Expose VI_B, Loop 1 English reconstruction.
% Scope: Counterexample 1.4.3 through Corollary 1.6.1 in full;
% Section 2 is excluded.
% Authority: Polo--Gille Exp6B-13oct24.pdf, component-local pages 7--8
% (combined-reader pages 348--349; original-margin printed pages 327--328).
% OCR and the Reinhold English Markdown are comparison witnesses only.

\medskip
\phantomsection\label{ctr:sga3-VIB-1.4.3}%
\noindent\textbf{\SGARefLink{ctr:sga3-VIB-1.4.3}{反例 1.4.3}.---}
\emph{设 \(k\) 为特征零的域, \(G\) 为常 \(k\)-群 \(\mathbb Z\),
\(H\) 为加法 \(k\)-群 \(\mathbb G_{a,k}\).
设 \(u:G\to H\) 为 \(k\)-群态射.
若 \(u\ne0\), 则 \(\ker u=0\), 但 \(u\) 不是闭浸入
(参见 \hyperref[ctr:sga3-VIB-1.2.2]{1.2.2}).}

下文将使用以下两个本应列入第~VI\(_A\)~篇的结果.

\medskip
\phantomsection\label{lem:sga3-VIB-1.5}%
\noindent\textbf{引理 1.5.---}
\emph{设 \(k\) 为域, \(G\) 为局部有限型的 \(k\)-群. 则:
\begin{enumerate}
\item[\textup{\SGAInlineRefTarget{lem:sga3-VIB-1.5:item:i}(i)}]
  \(G\) 的所有不可约分支都有相同的维数;
\item[\textup{\SGAInlineRefTarget{lem:sga3-VIB-1.5:item:ii}(ii)}]
  对每个 \(g\in G\),
  \[
    \dim_gG=\dim G.
  \]
\end{enumerate}}

\footnote{\label{note:sga3-VIB-17}编者注:
\textup{(ii)} \(\Rightarrow\) \textup{(i)} 是一般事实
(参见 EGA~0\(_{\mathrm{IV}}\), 14.1.6),
而 \textup{(i)} \(\Rightarrow\) \textup{(ii)} 来自以下事实:
若 \(X\) 为域上的有限型不可约概形, 则对 \(X\) 的每个非空开子集
\(U\), 都有 \(\dim X=\dim U\).
因而这里的关键是断言 \textup{\SGARefLink{lem:sga3-VIB-1.5:item:i}{(i)}},
它已在第~VI\(_A\)~篇 \SGARefLink{cor:sga3-VIA-2.4.1}{2.4.1}
的补充中建立. 因此相应修改了陈述及证明.}
断言 \textup{\SGARefLink{lem:sga3-VIB-1.5:item:i}{(i)}}
已在第~VI\(_A\)~篇, \SGARefLink{cor:sga3-VIA-2.4.1}{2.4.1} 中证明,
由此得到 \textup{\SGARefLink{lem:sga3-VIB-1.5:item:ii}{(ii)}}.
事实上, 设 \(g\in G\), \(C\) 为 \(G\) 中包含 \(g\) 的连通分支.
按照定义 (EGA~0\(_{\mathrm{IV}}\), 14.1.2),
\(\dim_gG\) 是整数 \(\dim U\) 的下确界,
其中 \(U\) 遍历 \(g\) 的开邻域. 因而对某个 \(U_0\), 有
\[
  \dim_gG=\dim U_0
\]
并且可假设它包含在 \(C\) 中
(因为当 \(V\subset U\) 时, \(\dim V\leq\dim U\)).
由于 \(C\) 不可约且在 \(k\) 上为有限型
(第~VI\(_A\)~篇, \SGARefLink{cor:sga3-VIA-2.4.1}{2.4.1}),
\[
  \dim U_0
    =\dim C
    =\operatorname{trdeg}_k\kappa(\xi),
\]
其中 \(\xi\) 是 \(C\) 的泛点 (EGA~IV\(_2\), 5.2.1).
因此
\[
  \dim_gG=\dim C=\dim G.
\]

\medskip
\phantomsection\label{prop:sga3-VIB-1.6}%
\noindent\textbf{命题 1.6.---}%
\footnote{\label{note:sga3-VIB-18}编者注:
在陈述中, 已将``沿单位截面''替换为``在单位截面的每一点处''.
在证明结尾, 还明确写出了所引的 EGA~IV\(_4\) 的结果.}
\emph{设 \(S\) 为特征零的概形, \(G\) 为 \(S\)-群概形,
在单位截面 \(\epsilon(S)\) 的每一点处在 \(S\) 上局部有限呈示.
概形 \(G\) 在单位截面的每一点处在 \(S\) 上平滑,
当且仅当称为 \(G\) 的单位截面的余法模的 \(\mathcal O_S\)-模}
\[
  \omega_{G/S}
    =\epsilon^*\bigl(\Omega^1_{G/S}\bigr),
\]
\emph{局部自由.}

回顾一下, 概形 \(S\) 称为特征零的,
是指对每个闭点 \(x\in S\), 域 \(\kappa(x)\) 的特征为零.

还回顾一下 (\SGARefLink{sga3:ii:subsection:4-11}{第~II~篇, 4.11}),
若 \(\pi:G\to S\) 为结构态射, 则
\[
  \Omega^1_{G/S}=\pi^*(\omega_{G/S}).
\]
因而, \(\mathcal O_S\)-模 \(\omega_{G/S}\) 局部自由,
等价于 \(\mathcal O_G\)-模 \(\Omega^1_{G/S}\) 局部自由.

若 \(\epsilon(S)\) 有一个在 \(S\) 上平滑的开邻域 \(U\),
则由 EGA~IV\(_4\), 17.2.3, \(\Omega^1_{U/S}\) 局部自由且为有限型,
从而
\[
  \omega_{G/S}=\epsilon^*(\Omega^1_{U/S}).
\]
也如此.

反之, 若 \(\omega_{G/S}\) 局部自由, 则
\[
  \Omega^1_{G/S}=\pi^*(\omega_{G/S}).
\]
也局部自由. 由于 \(S\) 特征为零, Jacobian 判别法
(EGA~IV\(_4\), 16.12.2) 表明 \(G\) 在 \(S\) 上微分平滑.
由 EGA~IV\(_4\), 17.12.5 可知,
\(G\) 在单位截面的每一点处在 \(S\) 上平滑.

\medskip
\phantomsection\label{cor:sga3-VIB-1.6.1}%
\noindent\textbf{推论 1.6.1 (Cartier).---}
\emph{设 \(k\) 为特征零的域.
每个在 \(k\) 上局部有限型的 \(k\)-群都在 \(k\) 上平滑.}

事实上, \(k\)-模 \(\omega_{G/k}\) 此时显然自由.
由 \SGARefLink{prop:sga3-VIB-1.6}{1.6}, \(G\) 在单位点 \(e\)
处在 \(k\) 上平滑, 从而由 \hyperref[cor:sga3-VIB-1.3.1]{1.3.1},
它在 \(k\) 上平滑.
